\newpage
\chapter{Factorization}

\emph{Factorization} is the process of writing a mathematical expression as a product of its factors. 
This is a fundamental technique in algebra, used to simplify expressions, solve equations, and analyze functions.

\subsection{Common Factor}

Factor out the greatest common divisor (GCD) of all terms.

\textbf{Example:}

\[
    6x^2 + 9x = 3x(2x + 3)
\]

\subsection{Difference of Squares}

\[
    a^2 - b^2 = (a - b)(a + b)
\]

\subsection{Sum of Cubes}

\[
    a^3 + b^3 = (a + b)(a^2 - ab + b^2)
\]

\subsection{Difference of Cubes}

\[
    a^3 - b^3 = (a - b)(a^2 + ab + b^2)
\]

\textbf{Example:}

\[
    x^3 + 8 = (x + 2)(x^2 - 2x + 4)
\]

\subsection{Trinomial: Special Case \texorpdfstring{\(x^2 + bx + c\)}{x² + bx + c}}

This is the case where \(a = 1\). Find two numbers whose product is \(c\) and whose sum is \(b\).

\textbf{Example:}

\[
    x^2 + 5x + 6 = (x + 2)(x + 3)
\]

\subsection{Trinomial: General Form \texorpdfstring{\(ax^2 + bx + c\)}{ax² + bx + c}}

As in the previous case, we are going to use a similar procedure, but we start by multiplying and dividing by \(\frac{a}{a}\), and then proceed with the rest. Look at the example:

\begin{align*}
3x^2 -5x - 2 &= \frac{3(3x^2 -5x - 2)}{3} \\
             &= \frac{{(3x)}^2 -5(3x) - 6}{3}\\ 
             &= \frac{(3x- 6)(3x+1)}{3} \\
             &= (x- 2)(3x +1)
\end{align*}

\subsection{Perfect Square Trinomial}

These follow the identities:

\[
    (a + b)^2 = a^2 + 2ab + b^2
\]

\[
    (a - b)^2 = a^2 - 2ab + b^2
\]

\textbf{Example:}

\[
    x^2 + 6x + 9 = {(x + 3)}^2
\]

\subsection{Substitution}

Substitute a more complex expression with a single variable, factor, then back-substitute.

\textbf{Example:}

Given \(x^4 + 2x^2 + 1\). Let \(y = x^2\), then

\[
    y^2 + 2y + 1 = (y + 1)^2  = (x^2 + 1)^2
\]

\subsection{Rationalization of Radicals}

Rationalizing removes radicals from the denominator.

\textbf{Example (Single Radical):}

\[
    \frac{1}{\sqrt{2}} = \frac{\sqrt{2}}{2}
\]

\textbf{Example (Binomial):}

\[
    \frac{1}{\sqrt{3} + 1} = \frac{\sqrt{3} - 1}{(\sqrt{3} + 1)(\sqrt{3} - 1)} = \frac{\sqrt{3} - 1}{2}
\]

\subsection{Horner’s Method}

This method is used to divide a polynomial by a binomial of the form \((x - r)\).

We can apply this method by following these steps:

\begin{enumerate}

    \item Write the coefficients of the polynomial \(P_1(x)\).
    \item Find the factors of the term with \(x^0\).
    \item Bring down the first coefficient.
    \item Multiply it by \(r\), and add to the next coefficient. Repeat until you reach the term with \(x^0\).
    \item If the final addition gives zero, then you have found a zero of the polynomial. Now, it can be written in the following form: \(P_2(x)(x - r)\), with \(P(x)\) being the new polynomial whose coefficients were computed during the process to get the zero.

\end{enumerate}

\textbf{Example:}

Given \(P_1(x) = x^3 - 6x^2 + 11x - 6\).

The factors of -6 are \(\{\pm 1, \pm 2, \pm 3, \pm 6\}\). We use \(r = 1\):

\[
  \begin{array}{r|rrrr}
  1 & 1 & -6 & 11 & -6 \\
    &   & 1 & -5 & 6 \\
  \hline
    & 1 & -5 & 6 & 0 \\
  \end{array}
  \implies P_2(x) = (x^2 - 5x + 6)(x - 1)
\]

Now, by repeating the process with \(r = 2\):

\[
  \begin{array}{r|rrr}
  2 & 1 & -5 & 6 \\
    &   &  2 & -6 \\
  \hline
    & 1 & -3 & 0 \\
  \end{array}
  \implies P_3(x) = (x - 3)(x - 2)(x - 1)
\]

The final result is \((x - 3)(x - 2)(x - 1)\).

\subsection{Long Division of Polynomials}

Another way of factorization is the so-called \emph{long division for polynomials}. This works exactly as the 
division algorithm for real numbers, but we use powers of \(x\) and the scalars of \(\field\) for the 
multiplication step.

\textbf{Example:}

\polylongdiv{x^3+x^2-1}{x-1}

with the result: \(x^2 + 2x + 2 + \frac{1}{x-1}\).
